%!TEX root = main.tex
\setcounter{chapter}{3}
\setcounter{section}{3}
\setcounter{theorem}{0}
\setcounter{equation}{0}

\lectureheader{161}{Calculus I}{Derivative rules}{\textit{Thomas' Calculus} \thesection, 3.5}

\begin{theorem}[Constant rule]
If $c\in\R$, then 
\begin{equation*}
\frac{\dee}{\dee x}c = 0.
\end{equation*}
\end{theorem}
\begin{theorem}[Power rule]
If $c\in\R$, then 
\begin{equation*}
\frac{\dee}{\dee x}x^c = cx^{c-1}.
\end{equation*}
\end{theorem}

\begin{example}
Differentiate the following with respect to $x$.
\begin{enumerate}[label=(\alph*)]
\item $y=x^{2/3}$
\vfill
\item $y=x^{\sqrt 2}$
\vfill
\item $y=x^{-4/3}$
\vfill
\item $y=\sqrt{x^{2+\pi}}$
\vfill
\end{enumerate}
\end{example}

\newpage

\begin{theorem}[Constant multiple rule]
If $u$ is a differentiable function of $x$ and $c\in\R$, then
\begin{equation*}
\frac{\dee}{\dee x}(cu) = c\frac{\dee u}{\dee x}.
\end{equation*}
\end{theorem}
\begin{proof}\,

\vspace{3in}
\end{proof}

\begin{theorem}[Sum rule]
If $u$ and $v$ are both differentiable functions of $x$, then
\begin{equation*}
\frac{\dee}{\dee x}(u+v) = \frac{\dee u}{\dee x} + \frac{\dee v}{\dee x}.
\end{equation*}
\end{theorem}
\begin{proof}\,

\vspace{3in}
\end{proof}

\newpage

\begin{example}
Find all places where the curve $y=x^4-2x^2+2$ has horizontal tangent lines.
\end{example}

\newpage

\begin{theorem}[Sine rule]\,
\begin{equation*}
\frac{\dee}{\dee x}\sin x = \cos x.
\end{equation*}
\end{theorem}
\begin{proof}\,

\vspace{3in}
\end{proof}

\begin{theorem}[Cosine rule]\,
\begin{equation*}
\frac{\dee}{\dee x}\cos x = -\sin x.
\end{equation*}
\end{theorem}
\begin{proof}\,

\vspace{3in}
\end{proof}

\newpage

\begin{theorem}[Product rule]
If $u$ and $v$ are differentiable functions of $x$, then
\begin{equation*}
\frac{\dee}{\dee x}(uv) = u\frac{\dee v}{\dee x} + v\frac{\dee u}{\dee x}.
\end{equation*}
\end{theorem}
\begin{proof}\,

\vspace{5in}
\end{proof}

\newpage

\begin{example}
Compute an equation for the tangent line to the graph of $y=x\sin x$ at $x=\pi/2$.
\end{example}

\newpage


\begin{theorem}[Quotient rule]
If $u$ and $v$ are differentiable functions of $x$, then
\begin{equation*}
\frac{\dee}{\dee x}\left(\frac{u}{v}\right) = \frac{v\frac{\dee u}{\dee x} - u\frac{\dee v}{\dee x}}{v^2}.
\end{equation*}
\end{theorem}
\begin{remark}
Students often like to memorize this one using the jingle 
\begin{center}
``low dee high minus high dee low over the square of what's below."
\end{center}
If you get it the wrong way around, the rhyme doesn't work.
\end{remark}
\begin{proof}\,

\vspace{5in}
\end{proof}

\newpage

\begin{example}
Compute the derivative of $y =\DS\frac{t^2-1}{t^3+1}$ with respect to $t$.
Then determine where the tangent lines to the graph have positive/negative/zero slope.
\end{example}

\newpage

\begin{example}
Compute the derivative of $f(x) = \frac{\cos x}{1-\sin x}$ and its domain.
\end{example}

\newpage

\begin{theorem}[Tangent, secant, cotangent, and cosecant rules]
\begin{align*}
\frac{\dee}{\dee x}\tan x &= \sec^2 x,\\
\frac{\dee}{\dee x}\sec x &=\sec x\tan x,\\
\frac{\dee}{\dee x}\cot x &=-\csc^2 x,\\
\frac{\dee}{\dee x}\csc x &=-\csc x\cot x.
\end{align*}
\end{theorem}
\begin{remark}\,
\begin{itemize}
\item After inspecting the domain of each derivative, we learn that the trigonometric functions are differentiable everywhere that they are defined.
\item Since differentiable functions are continuous, that proves that the trigonometric functions are continuous on their domains (a fact that we stated earlier without proof).
\end{itemize}
\end{remark}

\newpage

\begin{definition}[Higher derivatives]\,
\begin{itemize}
\item If $y$ is a function of $x$, then the \textbf{zeroth derivative of $y$ with respect to $x$} is the function $y^{(0)}=y$ itself.
\item If $n$ is a positive integer and $y$ is a function of $x$, then the \textbf{$n$th derivative of $y$ with respect to $x$} is the function
\begin{equation*}
y^{(n)} = \frac{\dee}{\dee x}y^{(n-1)}.
\end{equation*}
\end{itemize}
\end{definition}

\begin{remark}
If $y=f(x)$, we also use the notations
\begin{itemize}
\item $y^{(n)} = \frac{\dee^ny}{\dee x^n} = f^{(n)}(x) = \frac{\dee^n f}{\dee x^n}$,
\item $y^{(2)} = y'' = f''(x)$,
\item $y^{(3)} = y''' = f'''(x)$,
\end{itemize}
but using more than 3 primes is considered to be bad manners.
\end{remark}

\begin{example}
Compute $\DS\frac{\dee^2 y}{\dee t^2}$ if $y=\sec t$.
\end{example}
\vfill

\begin{example}
Compute all higher order derivatives of $y=x^3-3x^2+2$.
\end{example}
\vfill

